\section{Continuous model theory on spaces with finite nets}
\label{sec:continuous-models}

What numerical value should an alternating string of continuous quantifiers
have in first-order arithmetic? A single quantifier can be evaluated on finer
finite nets. Replacing quantifiers one at a time suggests an error proportional
to their number. Comparing the two whole evaluations by a relation gives a
sharper answer. The same atomic error controls every quantifier in the
formula. We first work entirely with finite rational data.

\paragraph{Language and interpretation.}
A formula has a specified number of free variables. An atom is a predicate
symbol applied to a tuple of these variables, with repetitions and permutations
allowed. A term operation can be included by giving the resulting atomic
predicate its own approximation contract. Use truth values in $[0,1]$,
rational constants, $1-u$, $\min$, $\max$, and the binders
$\inf_y,\sup_y$. More generally admit each fixed, rational-preserving PR
connective $C$ that maps $[0,1]^j$ into $[0,1]$ and is provably
$1$-Lipschitz for the maximum coordinate distance. These conditions include
averaging. The sup binder is also definable by $1-\inf_y(1-\varphi)$.
These are the continuous-logic quantifiers of Ben Yaacov, Berenstein, Henson
and Usvyatsov, \emph{Model theory for metric structures}, \S3
(\href{https://doi.org/10.1017/CBO9780511735219.011}{2008}).
The nonexpansive fragment gives the estimate below with coefficient one.
A connective of Lipschitz constant $K>1$ contributes a factor $K$ on
each path through its node.

\paragraph{Finite back-and-forth estimate.}
Let $U,V$ be two nonempty finite lists of possible quantified inputs,
each with its own rational atomic tables. The tables may also read free
parameters outside the lists.
Let $R$ be a relation between the two carriers such that each element of
$U$ has a relative in $V$ and each element of $V$ has a relative in $U$.
Suppose corresponding atomic values differ by at most $\delta$ whenever
all argument coordinates are $R$-related. Write $E_U\varphi$ for the
rational result of substituting minimum and maximum over $U$ for the
binders, with the indicated atomic tables. For coordinatewise related
assignments $a,b$,
\begin{equation}
  |E_U\varphi(a)-E_V\varphi(b)|\leq \delta .
  \label{eq:model-transfer}
\end{equation}

\begin{proof}
Induct externally on the formula. An atom has the assumed bound. For a
connective the maximum of the argument discrepancies is at most $\delta$,
and its nonexpansiveness gives the conclusion. For an infimum write
$f(u)=E_U\psi(a,u)$ and $g(v)=E_V\psi(b,v)$. Select the first
minimum $v_*$ of $g$ by rational comparisons and take any $u'$ related
to $v_*$. Extending the related assignments preserves their relation.
The induction hypothesis gives
\[
  \min_U f\ \leq f(u')\ \leq g(v_*)+\delta=\min_V g+\delta.
\]
Select the first minimum $u_*$ of $f$ and use the other direction of
coverage to get the reverse inequality. Each local choice disappears from
the conclusion. The supremum case follows by the complement formula.
In particular, no compatible matching of selected optimizers is needed.
\end{proof}

\paragraph{The metric presentation.}
Take a nonempty decidable set of number codes $D$ with a fixed PR rational
pseudometric $d$. For $h\geq0$ put $r_h=2^{-h}$. Supply a PR nonempty list
$U_h\subseteq D$ and prove coverage on $D$ within $r_h$. A computable
rounding map suffices. An existential proof of coverage in iPRA also suffices
in \eqref{eq:model-transfer}. For each atomic predicate in the fixed formula
supply rational samples $A_h(a)\in[0,1]$ and proofs that
\begin{equation}
 |A_h(a)-A_j(b)|\le Ld_\infty(a,b)+e_h+e_j,\qquad e_h=2^{-h}.
 \label{eq:model-atom}
\end{equation}
Here $d_\infty$ is the maximum coordinate distance, zero on empty tuples.
One natural $L$ can be chosen as the maximum of the finitely many atom
constants. Error-free atoms may use $e_h=0$ instead. All structure data are
fixed programs with numeric parameters, and each law above is a premise
with an iPRA proof in a concrete instance.

Write $E_h=E_{U_h}$, using the level-$h$ atomic tables. For any two
tuples of dense codes $a,b$ put $t=\max(d_\infty(a,b),r_h,r_j)$.
Relate two dense codes at distance at most $t$. Each element of $U_h$
can be rounded into $U_j$ within $r_j$, and conversely within $r_h$.
Every coordinate in a related tuple is at distance at most $t$, so
\eqref{eq:model-atom} bounds all atomic discrepancies. Thus
\begin{equation}
 |E_h\varphi(a)-E_j\varphi(b)|\leq Lt+e_h+e_j.
 \label{eq:model-name}
\end{equation}
Both levels simulate all quantifier moves with the same radius. This
explains why the bound is independent of quantifier rank.

A represented parameter tuple is a fixed PR family $a_h$ of dense codes
with $d_\infty(a_h,a_j)\le r_h+r_j$. Equation \eqref{eq:model-name}
bounds the value discrepancy by $(L+1)(r_h+r_j)$. Shifting the index by
$s$ with $2^s\ge L+1$ gives a regular real name
$k\mapsto E_{k+s}\varphi(a_{k+s})$. The same inequality with two
equivalent tuples shows equality of the resulting names: a requested
positive error is met by taking both radii and the cross-distance small.
The argument also compares distinct presentations. A two-sided relation
between their nets with atomic discrepancy $\delta_h\to0$ at a supplied
PR rate gives identical values for closed formulas of the fragment.
This supplies an approximate test for elementary agreement whenever
the required relations have been exhibited.

The Lipschitz hypothesis can be weakened to a common nondecreasing PR
modulus $\mu$: at input distance at most $2^{-\mu(v)}$ require atomic
discrepancy at most $2^{-v}+e_h+e_j$. Take
$H(n)=\max(n+3,1+\mu(n+3))$. For $m\ge n$, both net radii and the
distance between the tuple samples at levels $H(n),H(m)$ are at most
$2r_{H(n)}\le 2^{-\mu(n+3)}$. Use the same relation as before.
The resulting difference of formula values is at most
$3\cdot2^{-(n+3)}$. This gives a regular name even when no Lipschitz
constant exists. As before the index shift is independent of the number
of binders.

\paragraph{Why the represented binder has infimum semantics.}
Write $v_\varphi(a)$ for the name just constructed. Fix dense parameters
$a$ for clarity. Letting $j$ refine in \eqref{eq:model-name} shows
\[
    |v_\varphi(a)-E_h\varphi(a)|\le\eta_h,
    \qquad \eta_h=Lr_h+e_h .
\]
The passage to an inequality of names means that any additional dyadic
slack can be met by a sufficiently fine explicit index $j$.
For a fixed represented $a$ the same estimate holds for $E_h\varphi(a_h)$
because $d_\infty(a_h,a_j)\le r_h+r_j$.
Let $m=v_{\inf_y\psi}(a)$. Given any represented $x$, at level $h$
choose $u\in U_h$ within $r_h$ of $x_h$. The finite infimum is at most
$E_h\psi(a_h,u)$. Applying \eqref{eq:model-name} at the same level to
$u,x_h$ bounds this by $E_h\psi(a_h,x_h)+Lr_h+2e_h$.
Allowing $h$ to increase gives $m\le v_\psi(a,x)$ for every supplied
name $x$. For the converse approximation, the least minimizing sample
$w_h$ is a dense code with a constant name and
\[
 v_\psi(a,w_h)\le E_h\psi(a_h,w_h)+\eta_h
  = E_h(\inf_y\psi)(a_h)+\eta_h\le m+2\eta_h .
\]
Consequently every lower bound for all the $v_\psi(a,x)$ is at most $m$.
The witness function $h\mapsto w_h$ answers each requested optimization
accuracy. It gives a regular point name only if additional coherence is
proved. These two inequalities also establish agreement with the usual
infimum on the metric completion in any ambient semantics supporting that
completion: dense samples approximate every point and the atoms extend
uniformly. Internally each use substitutes particular arithmetic names and
compares rational observations.

\paragraph{A sentence with alternating quantifiers.}
On $[0,1]$, take rational dense points, the absolute-distance atom, and
the grid $U_h=\{i/N:0\le i\le N\}$ with $N=2^h$. The atomic constant
is $L=2$ under the tuple sup metric, by two triangle inequalities.
Consider the sentence
\[
 \inf_{a,b,c}\sup_x \min\{|x-a|,|x-b|,|x-c|\}.
\]
With $m\ge1$ centers the value on a grid of $N+1$ points is exactly $q/N$,
where $q$ is the least nonnegative integer satisfying
$m(2q+1)\ge N+1$. Indeed a single integer center covers at most $2q+1$
grid points at radius $q$. Conversely, partition the consecutive points
into blocks of length at most $2q+1$ and use the middle index of each
block, repeating centers if necessary. This proves the grid formula by a
finite count and a constructive cover. As $N$ grows, the inequalities
$N+1\le m(2q+1)$ and, for $q>0$, $m(2q-1)<N+1$ give
$q/N\to 1/(2m)$ with explicit error at most $1/N$.
The formula with three centers therefore has value $1/6$. The centers
$1/6,1/2,5/6$ realize it by their three closed intervals of radius
$1/6$. At $N=8,16,32$ the computed values are $1/8,3/16,5/32$
respectively. The program \texttt{model\_three\_centres.py} checks the
grid formula by independently enumerating all triples.

\paragraph{An infimum of zero without any computable minimizer.}
The repository's \href{https://github.com/VladimirReshetnikov/ProveIt/blob/5f3be02388442f16487d2dc5c2529961dd783d2b/Logic/PrimitiveRecursiveArithmetic/Derivations/Lean/IPRADerivations/Foundations/CoherenceObstruction.lean}{coherence obstruction} gives a PR tree
whose finite levels have PR selectors. It isolates a limitation of exact
witnesses even for a single Lipschitz atom on Cantor space. Number the
partial computable programs. At level $n$, admit a bit word $w$ of length
$n$ if for every $i<n$ the following implication holds. If program $i$
on input $i$ returns a bit $b$ within $n$ steps, $w(i)=1-b$.
Bounded simulation and the bounded scan are PR. Earlier outputs persist,
so prefixes of admitted words are admitted. At each level one can
explicitly fill every forced bit and set all remaining bits to zero.
A computable infinite path $p$ would have some program index $i$. After the computation of $p(i)$ finishes, a sufficiently long admitted
prefix would require $p(i)=1-p(i)$. This last statement is an external
computability argument about program indices.

Let $D$ consist of eventually-zero bit strings given by finite words with
zero padding. If two strings first differ at index $i$ their distance
is $2^{-i}$, and equal strings have distance zero. Length-$h$ words padded
by zeros form a PR $2^{-h}$-net. Define
\[
 A_h(x)=\sum_{n=1}^h 2^{-n}
       [\,x{\upharpoonright}n\text{ is not admitted}\,] .
\]
All summands are PR and lie in $[0,2^{-n}]$. The tail after $h$ is at
most $2^{-h}$. Inputs agreeing on the first $i$ bits have equal
summands through $i$, so the difference between the two limiting values
is at most $2^{-i}$. This proves \eqref{eq:model-atom} with $L=1$
by finite sums and their tails. An admitted word of length $h$ padded
by zeros has $A_h=0$ and limiting value at most $2^{-h}$. Hence the
represented infimum is identically zero as a real, with a PR approximate
minimizer at every precision. If a point had value zero, each prefix
would be admitted: the first failure at length $n$ would already give
a positive contribution $2^{-n}$. A computable regular name computes
each bit $i$ by taking a dense approximant at level $i+2$, since the
distance to the point is then strictly below $2^{-i}$. Thus no computable
regular name can realize the infimum. This converts the repository's
coherence difficulty into an exact boundary for a compact metric
structure with effective nets, uniformly small tails and an explicit
Lipschitz bound.

\paragraph{Arithmetic realization and Lean.}
For every fixed formula, compilation builds one numerical program
recursively on its finite syntax at the metalevel. Encode an assignment
as a finite coded tuple, extend by pairing, evaluate atoms by the supplied
PR symbols, and implement every infimum by a least-index argmin over
$0,\ldots,|U_h|-1$. A supremum uses the complemented cost.
At a fixed binder depth bounding every level by $M=|U_h|\ge1$, at most
$|\varphi|M^{\operatorname{depth}(\varphi)}$ tree nodes are evaluated
when the depth counts nested binders. This operation count is for a scan which stores its current best
cost and index, so that each child is evaluated once. The Lean
selector proves its functional specification and can recompute costs. The induction invariant asserts
that the chosen index is within the processed prefix and its cost is
at most each earlier cost. The latter is a bounded Boolean conjunction
of rational comparisons, hence admits quantifier-free induction. Clearing
positive denominators implements every comparison. The two local coverage
witnesses in \eqref{eq:model-transfer} are used by ordinary existential
elimination after the scan lemma. Neither an induction over a satisfaction
predicate inside iPRA nor a PR evaluator for all programs is needed.

\texttt{lean/ModelCase.lean} checks the selector and the complete formula
induction for tables in common integer units. Its definitions
\texttt{Corresponds}, \texttt{Related}, \texttt{value} and
\texttt{transfer} implement the two-sided relation, coordinatewise
assignments, evaluation and \eqref{eq:model-transfer}. For rational
tables with the displayed $\min,\max,1-u$ operations, multiply by
the product of the finitely many denominators (including the error
denominator) to obtain one integer unit. The generic Lean connective carries a proof of
integer nonexpansiveness. To interpret it on $[0,1]$ one must also ensure
that its outputs remain in the chosen scale $[0,S]$, as they do for the
included minimum and maximum. The rational-connective extension in the
manuscript uses exact rational codes and their proved nonexpansiveness.
For example, exact averaging may require a finer denominator at its
output. There is also a checked object proof for externally fixed finite nets in
\texttt{lean/ModelObject.lean}. For a rational threshold $q$, the upper
cut is the assertion that the computed value is at most $q$. At a
minimum or infimum this is the finite disjunction of child upper cuts,
and at a maximum or supremum it is their finite conjunction. These
equivalences follow by rational comparability and finite attainment.
Push $1-u$ to the atoms using the duality between min and max.
For a pair of related leaves use the source cut at $q$ and target
cut at $q+\delta$. The theorem \texttt{cutTransfer} derives the
corresponding implication for the whole fixed pair of tables in
\texttt{Prf PRAAx}. The rule \texttt{infCut} transports each disjunct forwards,
and \texttt{supCut} supplies a conjunct for each target input by
backward coverage. The reversed relation gives the other direction.
The code takes the atomic cut implications as derivations. Identifying
each composite cut with the numeric evaluation uses the rational-order
argument above. The constructor leaves
arithmetic parameters free, while it unrolls the finite net lists at the
Lean level. To uniformly internalize their numerical lengths, compile
each child to a \texttt{PR 2} symbol on a candidate index and a packed
environment. For the numerical inequality in common units, use the variable-horizon object theorem \texttt{UniformMinimum.transfer\_closed} from Section~\ref{sec:verification} in each direction. For an actual selected witness with independent rational denominators, use
\texttt{RationalArgmin.selectedF}, \texttt{selected\_spec} and
\texttt{selected\_minimal} from
the upstream \href{https://github.com/VladimirReshetnikov/ProveIt/blob/5f3be02388442f16487d2dc5c2529961dd783d2b/Logic/PrimitiveRecursiveArithmetic/Derivations/Lean/IPRADerivations/Analysis/RationalArgmin.lean}{rational argmin module}.
Instantiate the cross-table child theorem twice with the two minimizer
indices and locally opened covering indices. Additivity of rational order
and transitivity prove the two directed inequalities. Abstract over the
free numerical parameters. The metric and Cantor examples require their
specific PR net and table proofs in this translation. Equations
\eqref{eq:model-atom} and \eqref{eq:model-name} state the laws these
programs must export. The repository's report on
\href{https://github.com/VladimirReshetnikov/ProveIt/blob/5f3be02388442f16487d2dc5c2529961dd783d2b/Logic/PrimitiveRecursiveArithmetic/docs/proofs/foundations/ipra_embedding_framework_2026-11-15/paper/sections/08_existence_compactness.tex}{\emph{Finite nets, search, and exact objects from nested certificates}} treats a
single sampled residual. Here the relation estimate controls arbitrary
fixed nesting and explains exactly what the resulting quantifiers certify.
